% ==========================================
% LAMPIRAN B
% ==========================================
\cleardoublepage
\chapter{RINCIAN IMPLEMENTASI ALGORITMA \emph{AFFINITYENGINE} MODUL \emph{PEOPLE-TO-PROJECT}}
\label{chap:lampiran-b}

\renewcommand\thelstlisting{\thechapter.\arabic{lstlisting}}

\section{Alur Implementasi Algoritma}
\label{sec:b-alur-implementasi}

Implementasi algoritma dilakukan setelah sistem memperoleh data profil mahasiswa dan data kegiatan kolaboratif. Alur implementasi dimulai dari validasi data masukan, dilanjutkan dengan perhitungan skor kecocokan per atribut pada masing-masing sub-algoritma, penggabungan skor menjadi AffinityScore akhir, penerapan \emph{threshold} dan \emph{badge}, hingga perangkingan \emph{feed} rekomendasi. Ringkasan tahapan ini disajikan pada Tabel~\ref{tab:b1-alur}.

\begin{longtable}{p{0.06\textwidth} p{0.38\textwidth} p{0.44\textwidth}}
\caption{Ringkasan Alur Implementasi Algoritma \emph{AffinityEngine}} \label{tab:b1-alur} \\
\toprule
\textbf{Tahap} & \textbf{Nama Proses} & \textbf{Keluaran} \\
\midrule
\endfirsthead
\multicolumn{3}{c}{\parbox{0.88\textwidth}{\centering \tablename\ \thetable\ Ringkasan Alur Implementasi Algoritma (lanjutan)}} \\
\toprule
\textbf{Tahap} & \textbf{Nama Proses} & \textbf{Keluaran} \\
\midrule
\endhead
\midrule
\multicolumn{3}{r}{\textit{Bersambung ke halaman berikutnya.}} \\
\endfoot
\bottomrule
\endlastfoot
0 & Penyaringan kendala keras (\texttt{IsEligible}) & Status \emph{eligible}/tidak berdasarkan ketersediaan waktu, dijalankan sebelum perhitungan skor kecocokan \\
\midrule
1 & Validasi data masukan & Data profil mahasiswa dan data kegiatan kolaboratif siap digunakan \\
\midrule
2 & Pembentukan variabel atribut & \texttt{Ks}, \texttt{Kp} (\emph{skill}), \texttt{Ms}, \texttt{Mp} (minat), \texttt{expLevel} (pengalaman), \texttt{gayaKerja}, \texttt{preferensiPeran} \\
\midrule
3 & Perhitungan skor kecocokan Algoritma 1 (\emph{Demands-Abilities Fit}) & SkorSkill dan SkorPengalaman pada rentang $[0,1]$ \\
\midrule
4 & Perhitungan SkorAlgoritma1 & Skor gabungan berbobot 0,75/0,25 pada rentang $[0,1]$ \\
\midrule
5 & Perhitungan skor kecocokan Algoritma 2 (\emph{Needs-Supplies Fit}) & SkorMinat, SkorGayaKerja, dan SkorPeran pada rentang $[0,1]$ \\
\midrule
6 & Perhitungan SkorAlgoritma2 & Skor gabungan berbobot 0,611/0,194/0,194 pada rentang $[0,1]$ \\
\midrule
7 & Penggabungan AffinityScore akhir & AffinityScore pada rentang $[0,1]$ (kombinasi 50:50) \\
\midrule
8 & Penerapan \emph{threshold} 30\% & Status lolos/tidak lolos ambang batas rekomendasi (\texttt{TerapkanThreshold}) \\
\midrule
9 & Penentuan \emph{badge} tiga tingkat & Kategori \emph{badge} berdasarkan AffinityScore (\texttt{TentukanBadge}) \\
\midrule
10 & Perangkingan \emph{feed} & Daftar kegiatan terurut menurun berdasarkan AffinityScore \\
\end{longtable}

Catatan: penyaringan kendala keras (Tahap 0) merupakan mekanisme terpisah dari penerapan \emph{threshold} 30\% (Tahap 8). Tahap 0 memeriksa kelayakan dasar suatu kegiatan berdasarkan ketersediaan waktu mahasiswa sebelum skor kecocokan dihitung sama sekali (fungsi \texttt{IsEligible}), sedangkan Tahap 8 menyaring kegiatan yang telah memiliki AffinityScore namun berada di bawah ambang batas 30\% (fungsi \texttt{TerapkanThreshold}). Kedua mekanisme ini independen dan tidak boleh disatukan.

\section{Implementasi Perhitungan Skor Atribut}
\label{sec:b-implementasi-skor}

Perhitungan skor membedakan dua sub-algoritma: Algoritma 1 (\emph{Demands-Abilities Fit}) mencakup atribut \emph{skill} dan pengalaman, sedangkan Algoritma 2 (\emph{Needs-Supplies Fit}) mencakup minat, gaya kerja, dan preferensi peran. Listing~\ref{lst:b1-skor-skill}--\ref{lst:b8-breakdown} menyajikan \emph{pseudocode} seluruh fungsi implementasi secara berurutan sesuai alur pada Tabel~\ref{tab:b1-alur}.

\begin{minipage}{\textwidth}
\begin{lstlisting}[caption={{SkorSkill(profil, project)}}, label={lst:b1-skor-skill}, frame=single, captionpos=t, language=]
Function SkorSkill(profil, project)
  Input:  profil.skill        : himpunan tag skill milik mahasiswa
          project.skillDicari : himpunan tag skill dibutuhkan kegiatan
                                (gabungan skillDicari seluruh role)
  Output: skorSkill dalam [0,1]

  skorSkill <- Coverage(profil.skill, project.skillDicari)
  return skorSkill
End Function

Function Coverage(milik, dibutuhkan)
  if |dibutuhkan| = 0 then
    return 1  // tidak ada tuntutan -> skor penuh
  end if
  milikSet <- { lowercase(s) | s in milik }
  cocok    <- |{ d in dibutuhkan | lowercase(d) in milikSet }|
  return cocok / |dibutuhkan|
End Function
\end{lstlisting}
\end{minipage}

\begin{minipage}{\textwidth}
\begin{lstlisting}[caption={SkorPengalaman(profil, project)}, label={lst:b2-skor-pengalaman}, frame=single, captionpos=t, language=]
Function SkorPengalaman(profil, project)
  Input:  profil.pengalaman     : level pengalaman mahasiswa, ordinal 1..3
          project.pengalamanReq : level pengalaman dibutuhkan kegiatan, 1..3
  Output: skorPengalaman dalam [0,1]

  selisih        <- |profil.pengalaman - project.pengalamanReq|
  skorPengalaman <- 1 - (selisih / 2)  // 2 = selisih maksimum skala 1..3
  return skorPengalaman
End Function
\end{lstlisting}
\end{minipage}

\begin{minipage}{\textwidth}
\begin{lstlisting}[caption={SkorAlgoritma1(skorSkill, skorPengalaman)}, label={lst:b3-skor-algo1}, frame=single, captionpos=t, language=]
Function SkorAlgoritma1(skorSkill, skorPengalaman)
  Input:  skorSkill      : hasil SkorSkill, rentang [0,1]
          skorPengalaman : hasil SkorPengalaman, rentang [0,1]
  Output: skorAlgoritma1 dalam [0,1]

  BOBOT_SKILL      <- 0.75  // bobot ROC, Bab IV.3.2
  BOBOT_PENGALAMAN <- 0.25

  skorAlgoritma1 <- (BOBOT_SKILL * skorSkill)
                  + (BOBOT_PENGALAMAN * skorPengalaman)
  return skorAlgoritma1
End Function
\end{lstlisting}
\end{minipage}

\begin{minipage}{\textwidth}
\begin{lstlisting}[caption={SkorMinat, SkorGayaKerja, SkorPeran(profil, project)}, label={lst:b4-skor-minat-gk-peran}, frame=single, captionpos=t, language=]
Function SkorMinat(profil, project)
  Input:  profil.minatTag  : himpunan tag minat milik mahasiswa
          project.minatTag : himpunan tag minat kegiatan
  Output: skorMinat dalam [0,1]

  skorMinat <- Coverage(profil.minatTag, project.minatTag)  // Listing B.1
  return skorMinat
End Function

Function SkorGayaKerja(profil, project)
  Input:  profil.gayaKerja  : gaya kerja mahasiswa ("Terstruktur"|"Fleksibel")
          project.gayaKerja : preferensi gaya kerja kegiatan
  Output: skorGayaKerja dalam {0, 1}

  if lowercase(profil.gayaKerja) = lowercase(project.gayaKerja) then
    return 1
  else
    return 0
  end if
End Function

Function SkorPeran(profil, project)
  Input:  profil.preferensiPeran  : peran pilihan mahasiswa (taksonomi 3 kategori)
          project.rolesDibutuhkan : himpunan peran dibutuhkan kegiatan
  Output: skorPeran dalam {0, 0.5, 1}

  if |project.rolesDibutuhkan| = 0 then
    return 0.5  // kegiatan tak menetapkan peran spesifik -> netral
  end if
  rolesSet <- { lowercase(r) | r in project.rolesDibutuhkan }
  if lowercase(profil.preferensiPeran) in rolesSet then
    return 1
  else
    return 0
  end if
End Function
\end{lstlisting}
\end{minipage}

\begin{minipage}{\textwidth}
\begin{lstlisting}[caption={SkorAlgoritma2(skorMinat, skorGayaKerja, skorPeran)}, label={lst:b5-skor-algo2}, frame=single, captionpos=t, language=]
Function SkorAlgoritma2(skorMinat, skorGayaKerja, skorPeran)
  Input:  skorMinat     : hasil SkorMinat, rentang [0,1]
          skorGayaKerja : hasil SkorGayaKerja, dalam {0, 1}
          skorPeran     : hasil SkorPeran, dalam {0, 0.5, 1}
  Output: skorAlgoritma2 dalam [0,1]

  BOBOT_MINAT      <- 0.611  // bobot ROC, Bab IV.3.2
  BOBOT_GAYA_KERJA <- 0.194
  BOBOT_PERAN      <- 0.194

  skorAlgoritma2 <- (BOBOT_MINAT * skorMinat)
                  + (BOBOT_GAYA_KERJA * skorGayaKerja)
                  + (BOBOT_PERAN * skorPeran)
  return skorAlgoritma2
End Function
\end{lstlisting}
\end{minipage}

\begin{minipage}{\textwidth}
\begin{lstlisting}[caption={AffinityScoreAkhir(skorAlgoritma1, skorAlgoritma2)}, label={lst:b6-affinity-score}, frame=single, captionpos=t, language=]
Function AffinityScoreAkhir(skorAlgoritma1, skorAlgoritma2)
  Input:  skorAlgoritma1 : hasil SkorAlgoritma1, rentang [0,1]
          skorAlgoritma2 : hasil SkorAlgoritma2, rentang [0,1]
  Output: affinityScore dalam [0,1]

  BOBOT_ALGO <- 0.5  // kombinasi setara 50:50

  affinityScore <- (BOBOT_ALGO * skorAlgoritma1)
                 + (BOBOT_ALGO * skorAlgoritma2)
  return affinityScore
End Function
\end{lstlisting}
\end{minipage}

\begin{minipage}{\textwidth}
\begin{lstlisting}[caption={{IsEligible, TerapkanThreshold, TentukanBadge}}, label={lst:b7-gating}, frame=single, captionpos=t, language=]
Function IsEligible(profil, project)
  // Tahap 0 - kendala keras, dijalankan SEBELUM skor dihitung
  Input:  project.jadwalSlot       : himpunan slot waktu kegiatan
          profil.ketersediaanWaktu : himpunan slot ketersediaan mahasiswa
  Output: eligible dalam {TRUE, FALSE}

  if |project.jadwalSlot| = 0 then
    return TRUE  // kegiatan tak menetapkan jadwal -> lolos otomatis
  end if
  tersedia <- { lowercase(s) | s in profil.ketersediaanWaktu }
  return exists slot in project.jadwalSlot such that lowercase(slot) in tersedia
End Function

Function TerapkanThreshold(affinityScore)
  Input:  affinityScore : hasil AffinityScoreAkhir, rentang [0,1]
  Output: lolosFeed dalam {TRUE, FALSE}

  AMBANG_BATAS <- 0.3  // Bab IV.3.4, fixed rule
  return affinityScore >= AMBANG_BATAS
End Function

Function TentukanBadge(affinityScore)
  Input:  affinityScore : hasil AffinityScoreAkhir, rentang [0,1]
  Output: badge dalam {"hijau", "kuning", "abu", "none"}

  persen <- affinityScore * 100
  if persen >= 85 then return "hijau"
  if persen >= 60 then return "kuning"
  if persen >= 30 then return "abu"
  return "none"
End Function
\end{lstlisting}
\end{minipage}

\begin{minipage}{\textwidth}
\begin{lstlisting}[caption={{BreakdownKontribusi, PerangkinganFeed}}, label={lst:b8-breakdown}, frame=single, captionpos=t, language=]
Function BreakdownKontribusi(skorSkill, skorPengalaman, skorMinat,
                              skorGayaKerja, skorPeran)
  Input:  kelima skor atribut (masing-masing hasil Listing B.1-B.4)
  Output: breakdown : peta atribut -> {skor, kontribusi}

  BOBOT_ALGO <- 0.5
  kontribusi(bobotInternal, skor) <- BOBOT_ALGO * bobotInternal * skor

  breakdown["skill"]      <- {skor: skorSkill,      kontribusi: kontribusi(0.75,  skorSkill)}
  breakdown["pengalaman"] <- {skor: skorPengalaman, kontribusi: kontribusi(0.25,  skorPengalaman)}
  breakdown["minat"]      <- {skor: skorMinat,      kontribusi: kontribusi(0.611, skorMinat)}
  breakdown["gayaKerja"]  <- {skor: skorGayaKerja,  kontribusi: kontribusi(0.194, skorGayaKerja)}
  breakdown["peran"]      <- {skor: skorPeran,      kontribusi: kontribusi(0.194, skorPeran)}
  // total seluruh nilai kontribusi = affinityScore akhir
  return breakdown
End Function

Function PerangkinganFeed(daftarKegiatan)
  Input:  daftarKegiatan : kegiatan yang lolos IsEligible dan TerapkanThreshold
  Output: feed terurut menurun berdasarkan affinityScore

  sort daftarKegiatan by affinityScore descending
  return daftarKegiatan
End Function
\end{lstlisting}
\end{minipage}

\section{Contoh Penelusuran Perhitungan Skor}
\label{sec:b-contoh-penelusuran}

Bagian ini menyajikan satu contoh penelusuran implementasi untuk menunjukkan bagaimana nilai tiap atribut menghasilkan AffinityScore akhir. Contoh ini bersifat ilustratif dan digunakan untuk memperjelas alur implementasi, bukan sebagai hasil evaluasi akhir sistem.

\begin{longtable}{p{0.20\textwidth} p{0.33\textwidth} p{0.33\textwidth}}
\caption{Contoh Data Masukan Profil Mahasiswa dan Kebutuhan Kegiatan} \label{tab:b-contoh-input} \\
\toprule
\textbf{Atribut} & \textbf{Profil Mahasiswa} & \textbf{Kebutuhan Kegiatan} \\
\midrule
\endfirsthead
\multicolumn{3}{c}{\parbox{0.88\textwidth}{\centering \tablename\ \thetable\ Contoh Data Masukan (lanjutan)}} \\
\toprule
\textbf{Atribut} & \textbf{Profil Mahasiswa} & \textbf{Kebutuhan Kegiatan} \\
\midrule
\endhead
\midrule
\multicolumn{3}{r}{\textit{Bersambung ke halaman berikutnya.}} \\
\endfoot
\bottomrule
\endlastfoot
\emph{Skill}         & Python, Machine Learning, SQL      & Python, Machine Learning \\
\midrule
Minat                & AI, Data Science                   & AI, Data Science, Riset \\
\midrule
Gaya kerja           & Terstruktur                        & Terstruktur \\
\midrule
Pengalaman           & 2 (Menengah)                       & 2 (Menengah) \\
\midrule
Ketersediaan waktu   & Senin malam, Rabu sore             & Senin malam, Sabtu pagi \\
\midrule
Preferensi peran     & Developer                          & Developer, Designer \\
\end{longtable}

Sebelum skor dihitung, fungsi \texttt{IsEligible} memeriksa ketersediaan waktu: mahasiswa tersedia Senin malam dan Rabu sore, sedangkan kegiatan membutuhkan Senin malam atau Sabtu pagi. Karena terdapat irisan (Senin malam), mahasiswa dinyatakan \emph{eligible} dan perhitungan skor dilanjutkan.

\begin{longtable}{p{0.32\textwidth} p{0.13\textwidth} p{0.13\textwidth} p{0.25\textwidth}}
\caption{Hasil Perhitungan Skor AffinityEngine} \label{tab:b-contoh-skor} \\
\toprule
\textbf{Atribut} & \textbf{Skor} & \textbf{Bobot} & \textbf{Kontribusi} \\
\midrule
\endfirsthead
\multicolumn{4}{c}{\parbox{0.88\textwidth}{\centering \tablename\ \thetable\ Hasil Perhitungan Skor (lanjutan)}} \\
\toprule
\textbf{Atribut} & \textbf{Skor} & \textbf{Bobot} & \textbf{Kontribusi} \\
\midrule
\endhead
\multicolumn{4}{r}{\textit{Bersambung ke halaman berikutnya.}} \\
\endfoot
\bottomrule
\endlastfoot
\emph{Skill}                            & 1{,}000          & 0{,}75  & 0{,}750 \\
\midrule
Pengalaman                             & 1{,}000          & 0{,}25  & 0{,}250 \\
\midrule
\textbf{SkorAlgoritma1}                & \textbf{1{,}000} & {--}    & {--} \\
\midrule
\textit{× bobot Algoritma1}            & {--}             & 0{,}5   & \textbf{0{,}500} \\
\midrule
Minat                                  & 0{,}667          & 0{,}611 & 0{,}408 \\
\midrule
Gaya kerja                             & 1{,}000          & 0{,}194 & 0{,}194 \\
\midrule
Preferensi peran                       & 1{,}000          & 0{,}194 & 0{,}194 \\
\midrule
\textbf{SkorAlgoritma2}                & \textbf{0{,}796} & {--}    & {--} \\
\midrule
\textit{× bobot Algoritma2}            & {--}             & 0{,}5   & \textbf{0{,}398} \\
\midrule
\textbf{AffinityScore Total}           & {--}             & {--}    & \textbf{0{,}898} \\
\end{longtable}

Kolom Kontribusi pada baris atribut menunjukkan hasil Skor × Bobot internal atribut tersebut di dalam sub-algoritmanya. Jumlah kontribusi atribut dalam satu sub-algoritma menghasilkan SkorAlgoritma ($0{,}750 + 0{,}250 = 1{,}000$ untuk Algoritma 1; $0{,}408 + 0{,}194 + 0{,}194 = 0{,}796$ untuk Algoritma 2). Baris \textit{× bobot ke AffinityScore} menampilkan perkalian bobot setara 0,5 dari masing-masing SkorAlgoritma ke AffinityScore akhir (kombinasi 50:50 antara kedua algoritma). Pada contoh ini, \emph{skill} mahasiswa mencakup seluruh kebutuhan kegiatan sehingga skor \emph{coverage} bernilai penuh, minat hanya mencakup dua dari tiga tag kegiatan sehingga skornya $2/3 \approx 0{,}667$, sedangkan gaya kerja dan preferensi peran seluruhnya cocok. AffinityScore sebesar 0,898 (89,8\%) melampaui ambang batas 30\% dan masuk kategori \emph{badge} \texttt{"hijau"} ($\geq$ 85\%).
